% Part: lambda-calculus
% Chapter: introduction
% Section: beta

\documentclass[../../../include/open-logic-section]{subfiles}

\begin{document}

\olfileid{lam}{syn}{bet}
\olsection{$\beta$-తగ్గింపు}
\sourcecorrection{OLTELAMBETA-001}{మూల ఫైలు beta.texను వాక్యనిర్మాణ అధ్యాయం దిగుమతి చేస్తున్నా, దాని ఫైలు గుర్తింపులో అధ్యాయ సంకేతం intగా ఉంది. పాఠక విభాగపు శ్రేణితో సరిపడేలా తెలుగు లక్ష్యంలో దాన్ని synగా మార్చాం; స్థిర మూల ఫైలు, దాని పథం మారలేదు.}

$(\lambd[m][(\lambd[y][y]) m])$ను చూసినప్పుడు దానికి
$\lambd[m][m]$తో ఏదో సంబంధం ఉందని అనిపిస్తుంది: మొదటి పదాన్ని
“సరళీకరించిన” ఫలితం రెండవది కావాలి. \emph{$\beta$-తగ్గింపు}
అనే భావన ఈ అంతర్బోధకు ఆచారబద్ధ రూపం ఇస్తుంది.

\begin{defn}[$\beta$-సంకోచనం, $\bredone$] \ollabel{defn:betacontr}
  \emph{$\beta$-సంకోచనం} ($\bredone$) అనేది పదాలపై కింది
  షరతును తీర్చే కనిష్ఠ అనుకూల సంబంధం:
  \[
    (\lambd[x][N])Q \bredone \Subst{N}{Q}{x}
  \]
  $P \bredone Q$ అయితే $P$ పదం $Q$కు \emph{$\beta$-సంకోచితం}
  అయిందంటాం. $(\lambd[x][N])Q$ రూపంలోని పదాన్ని
  \emph{రెడెక్స్} అంటాం.
\end{defn}

\begin{prob} \ollabel{prob:def}
  \olref[lam][syn][alp]{defn:aconvone}లో బద్ధ చరం పేరు మార్పుకు
  ఇచ్చినట్లుగా, $\beta$-సంకోచనానికి సమానమైన ఆగమన నిర్వచనాలను
  పూర్తిగా రాయండి.
\end{prob}

\begin{defn}[$\beta$-తగ్గింపు, $\bred$] \ollabel{defn:betared}
  \emph{$\beta$-తగ్గింపు} ($\bred$) అనేది $\bredone$ను
  కలిగి ఉండే పదాలపై కనిష్ఠ స్వప్రావర్తక, సంక్రమణ సంబంధం.
  $P \bred Q$ అయితే $P$ పదం $Q$కు $\beta$-తగ్గిందంటాం.
\end{defn}

సందర్భం స్పష్టంగా ఉన్నప్పుడు $\bredone$కు బదులు $\redone$,
$\bred$కు బదులు $\red$ అని రాస్తాం.

అనౌపచారికంగా, $M \bred N$ అనేది $M$ నుంచి సున్నా లేదా
కొన్ని $\beta$-సంకోచన దశలతో $N$ను పొందగలిగినప్పుడే నిజం.

\begin{defn}[$\beta$-నియతమైనది]
ఇకపై $\beta$-సంకోచనం చేయలేని పదాన్ని
\emph{$\beta$-నియతమైనది} అంటాం.
\end{defn}

$M \bred N$ అయ్యి $N$ $\beta$-నియతమైతే, $N$ను $M$ యొక్క
\emph{నియత రూపం} అంటాం. ఒక పదపు నియత రూపం అనన్యమేనా అని
అడగవచ్చు; తరువాత చూడబోయినట్లుగా సమాధానం అవును.

కొన్ని ఉదాహరణలను పరిశీలిద్దాం.
\begin{enumerate}
\item మనకు
\begin{align*}
(\lambd[x][xxy]) \lambd[z][z] & \redone (\lambd[z][z])(\lambd[z][z]) y \\
& \redone (\lambd[z][z]) y \\
& \redone y
\end{align*}
\item ఒక పదాన్ని “సరళీకరించడం” దాన్ని మరింత సంక్లిష్టంగా చేయవచ్చు:
\begin{align*}
(\lambd[x][xxy])(\lambd[x][xxy]) & \redone (\lambd[x][xxy])(\lambd[x][xxy])y \\
& \redone (\lambd[x][xxy])(\lambd[x][xxy])yy \\
& \redone \dots
\end{align*}
\item అది పదాన్ని మార్చకుండా కూడా ఉంచవచ్చు:
\[
(\lambd[x][xx])(\lambd[x][xx]) \redone (\lambd[x][xx])(\lambd[x][xx])
\]
\item ఒక పదాన్ని ఒకటి కంటే ఎక్కువ విధాలుగా కూడా తగ్గించవచ్చు;
  ఉదాహరణకు,
\[
(\lambd[x][(\lambd[y][yx]) z]) v \redone (\lambd[y][yv]) z
\]
బాహ్య ప్రయోగాన్ని సంకోచిస్తే ఇది వస్తుంది; అలాగే
\[
(\lambd[x][(\lambd[y][yx]) z]) v \redone (\lambd[x][zx]) v
\]
అంతరంగ ప్రయోగాన్ని సంకోచిస్తే ఇది వస్తుంది. అయితే ఈ సందర్భంలో
రెండు ఫలిత పదాలూ అదే పదం~$zv$కు మరింతగా తగ్గుతాయి.
\end{enumerate}

చివరి ఉదాహరణలోని అంతిమ ఫలితం యాదృచ్ఛికం కాదు.
అది లాంబ్డా కలనశాస్త్రంలోని “చర్చ్--రోసర్ లక్షణం”
అనే లోతైన, ముఖ్యమైన లక్షణాన్ని చూపుతుంది.

\begin{digress}
  సాధారణంగా ఒక పదాన్ని $\beta$-తగ్గించడానికి ఒకటి కంటే
  ఎక్కువ మార్గాలుంటాయి. అందువల్ల అనేక తగ్గింపు వ్యూహాలు
  రూపొందాయి; వాటిలో ఎక్కువగా వాడేది \emph{సహజ వ్యూహం}.
  అది ఎల్లప్పుడూ \emph{అత్యంత ఎడమవైపు} ఉన్న రెడెక్స్‌ను
  సంకోచిస్తుంది. ఇక్కడ రెడెక్స్ స్థానం అంటే పదంలో అది
  మొదలయ్యే స్థానం. ఈ వ్యూహానికి ఉపయోగకరమైన లక్షణం ఉంది:
  ఏదైనా వ్యూహంతో పదాన్ని నియత రూపానికి తగ్గించగలిగితే,
  సహజ వ్యూహంతోనూ నియత రూపానికి తగ్గించగలం; తిరుగు దిశా
  సహజమే. ప్రత్యేకంగా చెప్పనప్పుడు ఇకపై సహజ వ్యూహాన్నే వాడతాం.
\end{digress}

\begin{defn}[$\beta$-తుల్యత, $\equal$]
  \emph{$\beta$-తుల్యత} ($\equal$) అనేది కింది విధంగా
  ఆగమన పద్ధతిలో నిర్వచించిన సంబంధం:
  \begin{enumerate}
  \item $M \equal M$.
  \item $M \equal N$ అయితే $N \equal M$.
  \item $M \equal N$, $N \equal O$ అయితే $M \equal O$.
  \item $M \equal N$ అయితే $PM \equal PN$.
  \item $M \equal N$ అయితే $MQ \equal NQ$.
  \item $M \equal N$ అయితే $\lambd[x][M] \equal \lambd[x][N]$.
  \item $(\lambd[x][N])Q \equal \Subst{N}{Q}{x}$.
  \end{enumerate}
\end{defn}

మొదటి మూడు నియమాలు దీన్ని తుల్యతా సంబంధంగా చేస్తాయి;
తరువాతి మూడు అనుకూల సంబంధంగా చేస్తాయి; చివరి నియమం
దీనిలో $\beta$-సంకోచనం కూడా ఉండేలా చేస్తుంది.

అనౌపచారికంగా, రెండు పదాలు $\beta$-తుల్యాలు కావడానికి,
ఒకదాన్ని సున్నా లేదా కొన్ని $\beta$-సంకోచన లేదా
తిరుగు $\beta$-సంకోచన దశలతో మరొకదిగా మార్చగలగడమే
అవసరమూ సరిపడేదీ. తిరుగు $\beta$-సంకోచనాన్ని ఇలా నిర్వచిస్తాం:
$M$ నుంచి $N$కు తిరుగు-$\beta$-సంకోచనం
అంటే $N$ నుంచి $M$కు $\beta$-సంకోచనం జరగడమే.

పై నియమాలకు మరిన్ని నియమాలను చేర్చి ఈ సంబంధాన్ని
విస్తరిస్తాం. అలా విస్తరించిన తుల్యతా సంబంధాన్ని
$\equal[X]$ అని సూచిస్తాం; ఇక్కడ $X$ చేర్చిన నియమం.

\end{document}
